% GRAFICO: due pannelli affiancati (una configurazione ciascuno), barre orizzontali, dati letti da un CSV,
% una sola legenda condivisa.
%
% Colonne di data.csv (x index): 0 e 1 = configurazione 1 (stato dell'arte, nuovo approccio); 2 e 3 = configurazione 2.
%
% VARIANTI
% - Altre configurazioni: il CSV ha anche le colonne 4–7 (configurazioni 3 e 4). Per mostrarle in una seconda slide
%   si copia questo file, si cambiano i titoli dei pannelli ("Configuration N") e gli x index (4, 5 e 6, 7).
% - Quattro pannelli (2x2): group size=2 by 2 e quattro \nextgroupplot. In una sola slide non entra (la riga di sotto esce
%   dalla pagina), quindi è adatto a una figura di un articolo. In quel caso, sotto ogni pannello si possono mettere
%   le lettere (a)-(d) con un ciclo dopo \end{groupplot}:
%       \foreach \letter/\target in {a/c1r1, b/c2r1}{ \node[below=0.25cm of my plots \target.south] {\small (\letter)}; }
%   dove cNrM è il pannello di colonna N e riga M del gruppo "my plots". Per la riga di sotto, che ha anche l'etichetta
%   dell'asse x, la distanza va portata a 0.8cm.
% - Legenda: è definita nel secondo pannello (legend to name=leg) e posizionata sopra i due pannelli con
%   \pgfplotslegendfromname. Lo stesso nome può essere definito una volta sola nel documento: se due slide includono
%   un grafico con questo codice, il secondo deve usare un nome diverso (per esempio leg2).
\pgfplotstableread[col sep=comma]{assets/tex/03_plot_quality/data.csv}\data

\begin{adjustbox}{max totalsize={.99\linewidth}{\textheight},center}
    \begin{tikzpicture}
        \pgfplotsset{every tick label/.append style={font=\normalsize}}
        \pgfplotsset{every axis label/.append style={font=\normalsize}}
        \begin{groupplot}[
                group style={
                    group name=my plots,
                    group size=2 by 1,
                    vertical sep=1.45cm, 
                    xticklabels at=all,
                    yticklabels at=edge left,
                    xlabels at=edge bottom,
                    ylabels at=edge left
                },
                xlabel={\certTimeSymbol\ (pm)},
                ylabel={Simulation},
                xticklabel style={
                    font=\small
                },
                y tick label style={
                    /pgf/number format/fixed, font=\small
                },
                %width=0.45*\textwidth,
                scaled y ticks=false,
                scaled x ticks=false,
                xbar, % to draw a horizontal histogram
                ytick=data,
                nodes near coords, % add data right of the bar
                enlarge x limits={upper, value=0.18}, % to have enough room for the nodes near coords, and just at the right-hand side
                %bar width=4pt,
                %xbar=4.5pt
                xmin=0,
        ]

            \nextgroupplot[title={\normalsize \textbf{Configuration 1}}]
            \addplot[draw=black, fill=gray!50, bar width=4.5pt] table[x index=0, y expr=\coordindex+1]{\data};
            \addplot[draw=black, fill=white, bar width=4.5pt] table[x index=1, y expr=\coordindex+1]{\data};

            \coordinate (top1) at (rel axis cs:0,1);% coordinate at top of the first plot

            \nextgroupplot[title={\normalsize \textbf{Configuration 2}},
                    legend style={
                        at={($(0,0)+(1cm,1cm)$)},
                        legend columns=2,
                        fill=none,
                        draw=none,
                        anchor=south,
                        inner sep=0.2em,
                        column sep=5pt,
                        align=center,
                        %yshift=1ex,
                    },
                    legend to name=leg,
            ]
            \addplot[draw=black, fill=gray!50, bar width=4.5pt] table[x index=2, y expr=\coordindex+1]{\data};
            \addlegendentry{\certTimeSymbol\ (\approachSOTA)}
            \addplot[draw=black, fill=white, bar width=4.5pt] table[x index=3, y expr=\coordindex+1]{\data};
            \addlegendentry{\certTimeSymbol\ (Our appr.)}

            \coordinate (top2) at (rel axis cs:1,1);% coordinate at top of the second plot

        \end{groupplot}

        % legend
        \coordinate (legend pos pre) at ($(top1)!.5!(top2)$);
        \coordinate (legend pos) at (legend pos pre |- current bounding box.north);
        \node[above] at (legend pos)
        {\pgfplotslegendfromname{leg}};

    \end{tikzpicture}
\end{adjustbox}